% ===================================================================
\section{Denotational and Operational Semantics}
\label{ch:semantics}
% ===================================================================

\subsection{Overview}

This chapter defines the semantic interpretation of FRFP terms. Two complementary
views are presented:

\begin{itemize}
  \item \textbf{Denotational semantics}: TDG terms denote morphisms in a
        Grothendieck construction over indexed categories.
  \item \textbf{Operational semantics}: pipeline execution is modeled as a
        relation on runs, connected to a stochastic trajectory space via the
        $\runToOmega$ bijection.
\end{itemize}

Both views contain significant structures made explicit by the Lean formalization
that were implicit or absent in the original published theory.

% -------------------------------------------------------------------
\subsection{The Grothendieck Construction}
% -------------------------------------------------------------------

FRFP uses a Grothendieck construction to model configurations --- the full state
of a pipeline execution including both categorical structure and index data.

\begin{definition}[Indexed category]
Define a contravariant functor:
\[
  \mathbf{E}_{\mathrm{idx}} : \mathbb{N}^{\mathrm{op}} \to \mathbf{Cat}
\]
mapping natural number indices $n$ to categories $\mathbf{E}_{\mathrm{idx}}(n)$
and index morphisms (decrements) to functors between them.
\end{definition}

\begin{definition}[Grothendieck object / configuration]
A configuration $\sigma$ is an object in the Grothendieck construction
$\int \mathbf{E}_{\mathrm{idx}}$:
\[
  \sigma = (n,\ X_n) \quad \text{where } n \in \mathbb{N},\
  X_n \in \mathrm{Ob}(\mathbf{E}_{\mathrm{idx}}(n))
\]
\leanref{Frfp.Core.Grothendieck.GrothendieckObject}
\end{definition}

\begin{definition}[Grothendieck morphism]
A morphism between configurations $(n, X_n) \to (m, X_m)$ consists of an index
morphism $n \to m$ and a morphism in the base category compatible with the
reindexing functor.
\leanref{Frfp.Core.Grothendieck.GrothendieckMorphism}
\end{definition}

% -------------------------------------------------------------------
\subsection{Lean Correction: Reduction Acts on Objects}
% -------------------------------------------------------------------

\begin{correction}[Lean Correction P-01]
The original published theory is ambiguous about whether reduction steps operate
on configurations, morphisms, or syntactic terms. The Lean formalization resolves
this: \textbf{reduction is a relation on configurations} (Grothendieck objects),
not on morphisms.
\end{correction}

\begin{definition}[Reduction step]
A reduction step is a relation
$\to_r \subseteq \mathrm{GrothendieckObject} \times \mathrm{GrothendieckObject}$.
\end{definition}

\begin{theorem}[Reduction induces morphism]
If $\sigma_1 \to_r \sigma_2$, then there exists a Grothendieck morphism $m$ with
$m.\mathrm{source} = \sigma_1$ and $m.\mathrm{target} = \sigma_2$:
\[
  \sigma_1 \to_r \sigma_2
  \Rightarrow
  \exists m : \mathrm{GrothendieckMorphism},\
  m.\mathrm{src} = \sigma_1 \wedge m.\mathrm{tgt} = \sigma_2
\]
\leanref{Frfp.Core.Grothendieck.reduction\_induces\_morphism}
\end{theorem}

% -------------------------------------------------------------------
\subsection{Semantic Functions}
% -------------------------------------------------------------------

\begin{definition}[Semantic evaluation]
The function $\Sem$ maps a configuration to an object in the kernel category:
\[
  \Sem : \mathrm{GrothendieckObject} \to \mathrm{Object}
\]
\leanref{Frfp.Core.TacitDependence.Sem}
\end{definition}

\begin{definition}[Correctness evaluation]
The function $\gamma$ maps a kernel object to a correctness judgment:
\[
  \gamma : \mathrm{Object} \to \mathrm{CorrectnessJudgment}
\]
\leanref{Frfp.Core.TacitDependence.gamma}
\end{definition}

\begin{definition}[Observable correctness]
The observable correctness function $\obs$ over runs:
\[
  \obs = \gamma \circ \Sem \circ \mathrm{initial}
       : \mathrm{Runs}(P) \to \mathrm{CorrectnessJudgment}
\]
\leanref{Frfp.Core.TacitDependence.obs}
\end{definition}

\begin{note}
$\Sem$, $\gamma$, $\obs$, and $\mathrm{CorrectnessJudgment}$ are all new
explicit types introduced by the Lean formalization. They were implicit in the
original paper's treatment of tacit-dependent correctness.
\end{note}

% -------------------------------------------------------------------
\subsection{Runs and Trajectories}
% -------------------------------------------------------------------

\begin{definition}[Run]
A run $\rho$ of program $P$ is a maximal execution trace. The type of runs of
$P$ is written $\mathrm{Runs}(P)$.
\end{definition}

\begin{definition}[Allowed stochastic trajectory]
The sample space $\Omega$ is the full space of stochastic trajectories. The
\emph{allowed} subspace $\Omega_{\mathrm{allowed}} \subseteq \Omega$ contains
exactly those trajectories consistent with a valid run of $P$.
\end{definition}

% -------------------------------------------------------------------
\subsection{The \texorpdfstring{$\runToOmega$}{runToOmega} Bijection}
% -------------------------------------------------------------------

\begin{correction}[Lean Correction P-02]
Remark A.104 in the original paper asserts a correspondence between runs and
allowed stochastic trajectories, but does not define it explicitly. The Lean
formalization requires a concrete injective map with a round-trip property.
\end{correction}

\begin{definition}[$\runToOmega$]
The run-to-trajectory correspondence:
\[
  \runToOmega : \{\rho \mid \mathrm{Runs}(P)\ \rho\} \to \Omega_{\mathrm{allowed}}
\]
with companion inverse:
\[
  \omegaToRun : \Omega_{\mathrm{allowed}} \to \{\rho \mid \mathrm{Runs}(P)\ \rho\}
\]
\leanref{Frfp.Core.OperationalSemantics.runToOmega}, \leanref{Frfp.Core.OperationalSemantics.omegaToRun}
\end{definition}

\begin{axiom_env}[$\runToOmega$ injectivity]
\[
  \runToOmega(\rho_1) = \runToOmega(\rho_2) \Rightarrow \rho_1 = \rho_2
\]
\leanref{Frfp.Core.OperationalSemantics.runToOmega_injective}
\end{axiom_env}

\begin{axiom_env}[$\runToOmega$ round-trip]
\[
  \omegaToRun(\runToOmega(\rho)) = \rho
\]
\leanref{Frfp.Core.OperationalSemantics.runToOmega_roundtrip}
\end{axiom_env}

\begin{note}
Injectivity and the round-trip property are axiomatized. A complete proof would
require formalizing the probability space construction. This is an open proof
obligation.

\textbf{Interpretation}: $\runToOmega$ is the formal bridge between the
operational semantics (discrete execution traces) and the probabilistic semantics
(measure on trajectory space). Without this bridge, probabilistic properties of
FRFP cannot be stated as theorems about runs.
\end{note}

% -------------------------------------------------------------------
\subsection{Denotational Semantics of Pipelines}
% -------------------------------------------------------------------

\begin{definition}[Pipeline denotation]
The denotational semantics $\llbracket \cdot \rrbracket$ maps a TDG term to its
morphism in $\Czero$:
\[
  \llbracket \varepsilon \rrbracket = \mathrm{id},\qquad
  \llbracket p \cdot t \rrbracket = \llbracket p \rrbracket \circ \llbracket t \rrbracket
\]
\end{definition}

\begin{theorem}[Denotational--operational agreement]
For any TDG term $t$ and its reduct $t'$, the denotations are equal:
\[
  t \to_c t' \Rightarrow \llbracket t \rrbracket = \llbracket t' \rrbracket
\]
\leanref{Frfp.Core.Semantics}
\end{theorem}

% -------------------------------------------------------------------
\subsection{Navigation Layer}
% -------------------------------------------------------------------

\begin{definition}[Navigation morphism]
A navigation morphism is an admissible path annotation on a TDG term, recording
which reachable states were visited during reduction. Navigation morphisms form a
groupoid.
\end{definition}

\begin{definition}[Navigation equivalence]
Two terms $t_1 \navequiv t_2$ if they are equal modulo navigation morphism
annotations.
\end{definition}

\begin{theorem}[Navigation equivalence is an equivalence relation]
\begin{itemize}
  \item Reflexivity: $t \navequiv t$.
  \item Symmetry: $t_1 \navequiv t_2 \Rightarrow t_2 \navequiv t_1$.
  \item Transitivity: $t_1 \navequiv t_2 \wedge t_2 \navequiv t_3
                      \Rightarrow t_1 \navequiv t_3$.
\end{itemize}
\leanref{Frfp.Core.SemanticCorrectness.nav\_equivalence\_is\_equivalence}
\end{theorem}

% -------------------------------------------------------------------
\subsection{Immutability of Pipelines}
% -------------------------------------------------------------------

\begin{theorem}[Pipeline immutability]
Executing a pipeline does not modify its structure. A pipeline's source, target,
and primitive are invariant under execution:
\[
  \mathrm{execute}(\Pi).\mathrm{structure} = \Pi.\mathrm{structure}
\]
\leanref{Frfp.Core.ImmutabilityProof}
\end{theorem}

This theorem captures a key FRFP invariant: pipelines are static specifications.
Execution produces outputs but does not rewrite the pipeline specification itself.
